\begin{proposition}\label{prop:octic-C55}
There is no point of order $55$ on an elliptic curve over a number field
of degree $8$.
\end{proposition}
\begin{proof}
Let $X=X_1(55)$ and let $D$ be the rational effective divisor of degree
$8$ arising from such a point. Its support is noncuspidal. The complete
newform decomposition has dimension $81$, and every constituent has
proved analytic rank zero. Consequently $J_X(\mathbb Q)$ is finite.
The input and the dimension check are recorded in
\texttt{C55\_rank\_certificate.json}.

At $2$, the only noncuspidal closed points of degrees at most $8$ have
degree $7$. There are $20$ of them, forming one diamond orbit. The
cuspidal degree signature is
\[
 1^{20},\quad4^5,\quad10^2,\quad20.
\]
All $20$ rational cusps lift to rational cusps over $\mathbb Q$.
Suppose first that $\overline D$ is not entirely cuspidal. Then it is
$P+c$, where $\deg P=7$ and $c$ is a rational cusp. Put
$A_3=T_3-1-3\langle3\rangle$. Eichler--Shimura and injectivity of torsion
reduction at the odd prime $3$ show that $A_3J_X(\mathbb Q)=0$.
Moreover, $A_3c=0$ as a divisor: on a rational cusp, the Hecke
correspondence consists of one copy of the cusp and three copies of
its diamond translate. Thus $A_3P$ must be principal. The fresh
calculation \texttt{C55\_inputs/N55\_D8\_NC7\_q3.log} proves that this is
false for a representative of the unique diamond orbit. Since $A_3$
commutes with diamond operators, this excludes this case.

Now suppose that $\overline D$ is entirely cuspidal. Let
\[
 H=\langle\langle16\rangle\rangle,\qquad Y=X/H.
\]
The group $H$ has order $5$ and acts freely. This remains true in
characteristic $2$: no nontrivial elliptic-curve automorphism has order
$5$, and the primitive-column calculation checks freeness at every
cusp. Hence $Y$ has good reduction at $2$ and genus $17$.
Its Jacobian has rank zero, since pullback embeds it, up to isogeny,
in $J_X$.

The map $Y\longrightarrow X_0(55)$ is an etale cyclic cover of degree
$4$, with generator induced by $\langle2\rangle$. There are no
elliptic stabilizers: an automorphism of order $4$ would require a
square root of $-1$ modulo $11$, and an automorphism of order $3$
would require a nontrivial cube root of unity modulo $5$. The
primitive-column calculation checks the cusps, including the special
fibre at $2$. In the order of cusp widths $55,11,5,1$, the Frobenius
elements in this cyclic group are
\[
 0,\quad1,\quad2,\quad3.
\]
The resulting cusp signature on $Y_{\mathbb F_2}$ is
$1^4,2^2,4^2$.

For completeness, the special-fibre model is identified independently
of the class-field constructor. The exact plane model and its
$j$-function are retained in \texttt{C55\_builtin\_model.m}; its
normalization has genus $5$, and its degree-$72$ $j$-map has the four
rational poles of widths $1,5,11,55$. The two modular-polynomial
identities for $5$ and $11$ are checked exactly. Since the normalization
of $X_0(5)\times_{X(1)}X_0(11)$ is the connected degree-$72$ cover
$X_0(55)$, these identities and the degree check also identify the
base function field. Its Picard group is
$\mathbb Z/5\mathbb Z\times\mathbb Z/20\mathbb Z\times\mathbb Z$.
There is exactly one character to $\mathbb Z/4\mathbb Z$ with the
specified cusp Frobenius elements. The corresponding length-two Witt
equation is converted to an explicit quartic, with exact inverse
maps and an automorphism whose square sends the quartic generator
to itself plus $1$. Actual Witt Frobenius traces on two places
generate the Picard group modulo $4$ and identify the character,
up to inversion. The quartic has genus $17$ and a split rational
cusp. Riemann--Hurwitz therefore proves that the cover is everywhere
unramified and geometrically connected. These checks identify it
with $Y_{\mathbb F_2}$ by global class field theory.

Let $\pi:X\longrightarrow Y$. The divisor $\pi_*\overline D$ is
supported on the four rational cusps of width $55$ and the one
quartic cusp of width $11$. Indeed, the other cusps on $Y$ are images
of source cusps of degrees $10$ or $20$, which cannot occur in a
divisor of degree $8$. This restriction, including all residue-degree
multiplicities, is checked using primitive-column cusp labels in
\texttt{cusps\_order5\_quotients.py}. There are precisely
\[
 [z^8](1-z)^{-4}(1-z^4)^{-1}=201
\]
possible image divisors.

On the verified genus-$17$ model put
$U=(\mathcal C-1)H^0(Y_{\mathbb F_2},\Omega^1)$.
Its dimension is $15$. Exact differential-space computations give
the $34$ local restriction kernels needed for the five admissible
closed cusps. The restriction of $U$ to each of the $201$ divisors
has rank $8$. The quotient Cartier criterion, applied to
$\pi_*D$, excludes all these reductions. Its hypotheses hold:
$Y$ has rank zero, a rational cusp, good reduction at $2$, and an
etale cyclotomic cuspidal boundary; the Hecke and diamond operators
descend from $X$. The kernel equalities, the Cartier projections,
and all $201$ ranks are replayed independently over $\mathbb F_2$.
This excludes the entirely cuspidal case and proves the proposition.
\end{proof}
